% 02-reading.tex: how to read the map.
\section{How to read this map}
\label{sec:reading}

Polaris is not one enormous table or one enormous server. It is a very small credential core with
deliberately bounded structures arranged around it, each of which constrains, protects, observes,
verifies, records or limits something about the structure inside it. The paper reads outward from
the core, ring by ring (Figure~\ref{fig:core-to-shell}), and the reader should be able to zoom:
from the person, to the credential, to the cryptography that authenticates it, to the schema rules
that bound it, to the authority that issued it, to the parties that rely on it, to the federation of
authorities, to the witnesses that watch them, to institutional exchange, to the operations that
keep it alive, to the society it sits inside, to the environment of software agents it now serves,
and finally to the outside: the software written by strangers whose acceptance or refusal is the only
evidence the repository did not produce about itself.

\subsection{The questions asked of every layer}

For each major component the paper answers the same questions, in a card of the form used
throughout. What is this? Why does it exist? What problem does it solve? What does it trust? What
does it refuse to trust? What can it see? What is it deliberately prevented from seeing? What happens
if it fails? What part of Polaris checks it? Is it needed now, or is it preparation for a future
environment? How does it connect to the layer inside it and the layer outside it? A reader who can
answer those questions for every ring has the mental model this paper exists to build.

\subsection{The town}

\analogy{The credential record is the register kept in the town hall: one entry per person, the
claim that everything else is organised around. The wall is the schema: triggers, check constraints
and unique indexes that decide which states may exist at all, and which bind every visitor, every
official and every restore from backup alike. The gates are the endpoints through which a claim
enters or leaves: issuance, verification, login, exchange. The watchtowers are the transparency logs,
and the witnesses stand outside the wall, cosigning what the towers publish and shouting when two
towers disagree. Other towns are other authorities, each with its own hall and its own key, joined by
roads that run in one direction and for one purpose. Above everything stand laws that the town itself
cannot repeal by a vote of its officials. And beyond the walls are travellers from towns nobody here
built, who owe the town nothing; what they find at its gate is the last thing this map shows.}

The map of the town, with the technical name beside every civic one, is Figure~\ref{fig:town-map}
in Section~\ref{sec:walls}. The analogy will recur in a shaded box like the one above, and every
recurrence is followed by the mechanism it stands for. Each ring also opens with one line in italics:
the paper's single idea, restated for that ring, that the ring inside it was correct and still not
enough.

\figfile{core-to-shell}{The concentric map of Polaris. The core is one credential record and its
signature. Each ring names the actual components on it, and the reading order of the paper follows
the rings outward. The outermost ring is dashed because most of it is not code: it is what parties
outside the repository have done with the rings inside it, and at \code{1.0.0-rc.7} that is one
wallet, one hosted conformance run and one published test corpus.}{fig:core-to-shell}

\subsection{Trust and refusal}

The most useful habit for reading Polaris is to ask, of every layer, what it refuses to take on
faith from the layer below. The schema refuses to trust the application. The detached verifier
refuses to trust the issuer's server. The second witness refuses to trust the first implementation.
The relying party refuses to trust a transitive chain of authorities. The witness refuses to trust
the log's own head. The check layer refuses to trust a check that cannot detect its own violation.
The mutation drills refuse to trust a test that has never been shown to fail. And the scoreboard
refuses to count anything the author ran against the author's own code.
The conclusion draws the refusals as a ladder (Figure~\ref{fig:layer-ladder}); Sections~\ref{sec:core} through
\ref{sec:operations} explain each one.
